\section{Further results}
\label{sec:further}

\subsection{Aspect ratio of the kite}
\label{sec:res_aspect}

The kite of every result above keeps the base aspect ratio of 3, its span squared over its area (\cref{sec:kite}). \Cref{fig:aspect} asks one what-if on three ships: the same kite area at a higher aspect ratio, with the launch gear, the tether, the canopy and everything else priced as at the base value. It shows what a better aspect ratio is worth if such a kite can be handled at the same cost, not that the kite should fly at that value. On that assumption the saving of the base kite on the \SI{140}{m} ship goes from \tbd{13.0}\,\% at aspect ratio 3 to \tbd{20.1}\,\% at 5 and \tbd{21.1}\,\% at 8. The canopy mass per area and the force limit, the largest tether force allowed, stay as they are, and the polars are those of a design shape that does not deform. \tbd{Draft numbers, at the canopy price goal, with the kite held at \SIlist{3200;2828;6889}{m^2} on the three ships; the final run replaces them.}

\begin{figure}[tbp]
\centering
\plot[\columnwidth]{aspect}
\caption{What-if: cost saving of the kite at the same area and the same price against its aspect ratio, on the \SI{140}{m} e-ammonia ship to Guadeloupe, a Kamsarmax on VLSFO to Tubarão and a Neo-Panamax container ship on e-ammonia to New York.}
\label{fig:aspect}
\end{figure}

\subsection{Force limit of the kite}
\label{sec:res_force}

Every kite of the main results is built for \SI{0.5}{kN/m^2}. A kite built for more force has a thicker tether and bridle, a dearer winch and heavier, dearer cloth (\cref{sec:kite}), and it depowers in stronger winds only. \Cref{fig:force_rating} routes the three ships of \cref{fig:aspect} with kites built for \SIrange{0.3}{1}{kN/m^2}, each at its best of five kite areas a factor $\sqrt2$ apart. A lower limit saves more on all three ships, and the best area grows as the limit falls. A kite built for more force is heavier, so it needs more wind to fly and is up for less of the passage. The share of time aloft falls from about 0.57 to 0.39 on the Neo-Panamax container ship between 0.3 and \SI{1}{kN/m^2}.

Doubling the limit to \SI{1}{kN/m^2} costs 5.4, 2.3 and 3.2 points on the \SI{140}{m} ship, the Kamsarmax and the Neo-Panamax, far beyond the scatter of the pattern search, about 0.5 points on a difference between two savings. Lowering it to \SI{0.3}{kN/m^2} gains 0.8, 1.0 and 0.9 points, about twice that scatter and of the same sign on every ship. On the Kamsarmax and the Neo-Panamax the best kite at that limit is the largest searched, so the gain there is a lower bound. Most of the loss at a higher limit comes from the tether, bridle and winch, which keep the kite down in light wind. The cloth is the smaller part. On the same areas, cloth fixed at \SI{200}{g/m^2} whatever the limit, with the seams and reinforcements still growing with it, changes the saving at \SI{1}{kN/m^2} by \tbd{+2.1}, \tbd{+0.1} and \tbd{-0.2} points. Cloth in proportion to the limit, a constant share of its strength, costs \tbd{3.0}, \tbd{1.3} and \tbd{1.5} points more than the decay law, the loss of cloth strength to sunlight and wear that sets the canopy's life (\cref{eq:uv_life}). The \SI{140}{m} ship gains most from lighter cloth because its kite is held down by its mass. \tbd{These savings were found with a control unit of fixed mass; at \SI{0.22}{kg} per kN of force limit it weighs \SI{88}{kg} on a kite of \SI{800}{m^2} built for \SI{0.5}{kN/m^2} and \SI{176}{kg} built for \SI{1}{kN/m^2}, so the final run will charge a higher limit somewhat more.} The bounds of \SIrange{0.4}{1}{kN/m^2} on the force limit in \cref{tab:inputs} carry the same cloth, seams and reinforcements.

We keep \SI{0.5}{kN/m^2} for every other result. The canopy goals of life and cloth weight are stated at that limit, and the gain from a lower one is within twice the scatter. The kite's savings in \cref{sec:results} may therefore be about a point short of what a kite built for a lower force would reach. \tbd{Draft numbers; the final run replaces them.}

\begin{figure}[tbp]
\centering
\plot[\columnwidth]{force_rating}
\caption{Cost saving of the kite at its best area against the force limit it is built for, on the ships of \cref{fig:aspect}, with cloth that keeps its life under the decay law (solid), in proportion to the limit (dashed) and fixed at \SI{200}{g/m^2} (dotted). The dashed and dotted lines take the best of three areas around the ship's own kite, the solid line of five.}
\label{fig:force_rating}
\end{figure}

\subsection{Kite and rotors together}
\label{sec:res_combination}

Kite and rotors pull best at different wind angles (\cref{fig:rig_polar}), and a kite needs only a launch mast at the forward end of the deck, so it can be added to a row of rotors. \Cref{fig:combination} compares the best pair of kite and rotors with the best of each rig alone, from kites of \SIlist{400;800;1600;3200}{m^2} (the pair up to \SI{1600}{m^2}) and rotors filling a fifth, half and all of the deck left behind the launch mast. On e-ammonia the pair saves \tbd{2.1} points more than the better single rig on both the \SI{140}{m} ship and the Kamsarmax, up to \tbd{4.0} points on the best route. On VLSFO it saves \tbd{0.2} points more and \tbd{0.04} less. It keeps \tbd{about 60}\,\% of the sum of the two single-rig savings on e-ammonia and \tbd{about half} on VLSFO. \tbd{This draft ran before the sign of the kite's side force was corrected, which made a kite over rotors cancel side force that it adds; the final run replaces these numbers and the reading below.} The likely reason, which we did not separate from the data, is that the two rigs saturate. In a moderate wind either alone covers most of the propulsive power the ship needs at the speed the router picks. The kite adds least where the rotors leave most to gain, since it does not fly in light wind and gives little close-hauled. Where both pull, their side forces add up in the hull's balance, and the extra leeway and rudder drag take back part of the thrust. The kite's side force at the launch mast, ahead of the point along the hull at which the hull's own side force acts, offsets some of the rotors' yaw moment, their turning moment about the vertical axis, close-hauled. Each rig keeps its full capital and annual cost, so the pair pays twice for gains that overlap.

\begin{figure*}[tbp]
\centering
\plot[\textwidth]{combination}
\caption{Kite and rotors together on the \SI{140}{m} ship and the Kamsarmax bulk carrier with a VLSFO and an e-ammonia drive: (a) cost saving against the ship without a rig of the best kite of \SIlist{400;800;1600;3200}{m^2}, the best rotors at a fifth, half and all of the deck capacity left behind the kite's launch mast, and the best pair of a kite of \SIlist{400;800;1600}{m^2} and those rotors flown together; (b) the saving of that pair over the better of the best kite and the best rotors alone. Bars are the mean over the six routes, whiskers the range over them.}
\label{fig:combination}
\end{figure*}

\subsection{Kite savings at a held speed}
\label{sec:further_speed}

\Cref{tab:speed_size} gives the values behind \cref{fig:speed_size} (\cref{sec:res_speed}): for each ship, fuel and held speed, the smallest kite, the kite that saves the most fuel and the priced kite that saves the most cost. We hold the \SI{140}{m} ship at 8 to \SI{11}{kn} on VLSFO and e-ammonia, with kites of 400 to \SI{6400}{m^2}, and the feeder container ship at 13, 16 and \SI{19}{kn} on VLSFO, with kites of 800 to \SI{12800}{m^2}, on the routes to Guadeloupe and New York. CE Delft, by comparison, fit one \SI{400}{m^2} kite to every ship \citep{cedelft2016}. On the \SI{140}{m} ship the kite that saves the most fuel grows from \SI{3200}{m^2} to \SI{6400}{m^2} between 8 and \SI{11}{kn}. At a held speed a kite cannot save more than the fuel spent on propulsion, which is least at the lowest speed, so the largest kite gains most with speed. Priced on e-ammonia the cheapest kite saves \tbd{5.5}\,\% of the cost at \SI{8}{kn} and \tbd{9.5}\,\% at \SI{11}{kn}. On VLSFO kites of \SI{3200}{m^2} and more cost more than they save at all of these speeds. On the feeder the fuel grows by \tbd{90}\,\% from 13 to \SI{19}{kn} and the energy its smallest kite saves falls by \tbd{25}\,\%. The \SI{6400}{m^2} kite saves the most at every speed, \tbd{21.5}\,\% at \SI{13}{kn} and \tbd{19.2}\,\% at \SI{19}{kn}, a loss of \tbd{2.3} points. Free of running cost it is up \tbd{35 to 38}\,\% of the time at every speed and the \SI{12800}{m^2} kite \tbd{20 to 23}\,\%. The apparent wind moving ahead costs no time aloft in these ranges, while \citet{podeur2016} report less flight time at \SI{19}{kn}. Its cost is in the thrust per unit of kite force, which a larger kite makes up for as far as these limits allow.

\begin{table*}[tbp]
\centering
\caption{Kite savings at a held ship speed, on both routes: fuel of the ship without a rig per round trip, fuel saved by the smallest kite and by the kite that saves the most fuel, both free of running cost, the priced kite that saves the most cost, and the share of the time the largest kite is up. \tbd{Draft run.}}
\label{tab:speed_size}
\InputIfFileExists{sections/speed_size_table}{}{\tbd{Speed and size table pending: run the speed\_size\_table step.}}
\end{table*}
